\section*{Chapter \ref{ch:b2:behavioural-evidence} --- Behavioural Evidence: Where Rational Choice Breaks}

\begin{solution}{exo:b2:behavioural-evidence:1}
% code: exercises/solutions.py ex_replication
Fifteen of the 18 dots lie below the diagonal, and three below the dotted line (an effect of the opposite
sign). The \omterm{def:b2:behavioural-evidence:replication-rate}{replication rate} is $11/18 = 61.1$ per cent. A real effect is detected only with a
probability equal to the power of the test, and when original estimates passed a significance filter,
the replications were powered for effects larger than the true ones
(\cref{prop:b2:behavioural-evidence:filter}). Both lower the rate even if every effect is real.
\end{solution}

\begin{solution}{exo:b2:behavioural-evidence:2}
% code: exercises/solutions.py ex_mixed
$V = \tfrac12 \times 150 - \tfrac12 \times 2 \times 100 = 75 - 100 = -25$: she refuses. She accepts when
$\tfrac12 z_+ \ge \tfrac12 \times 2 \times 100$, that is from $z_+ = \dal{200}$.
\end{solution}

\begin{solution}{exo:b2:behavioural-evidence:3}
% code: exercises/solutions.py ex_endowment
An owner gives up the mug, a loss valued at $\kappa u = 10$: $\mathrm{WTA} = \dal{10}$. If money is
neutral, a buyer pays up to the gain $u$: $\mathrm{WTP} = \dal{5}$, a ratio of 2. If paying is also a
loss, she pays $P$ with $\kappa P = u$: $\mathrm{WTP} = \dal{2.5}$, a ratio of $4 = \kappa^2$.
\end{solution}

\begin{solution}{exo:b2:behavioural-evidence:4}
For the gamble: the same person accepts the gamble when it is part of a series she evaluates as a whole,
or turns down gambles with no loss at all (only a gain or nothing), which \omterm{def:b2:behavioural-evidence:loss-aversion}{loss aversion} cannot explain.
More directly, a person who accepts 50--50 gambles with $z_+$ only slightly above $z_-$ shows $\kappa$
close to one. For the WTA--WTP gap: the gap vanishes when the good is a token with a known redemption
value, or when the owners are dealers who hold the good for resale; a gap that is the same for tokens
and goods points to a cost of trading or a misunderstanding of the procedure, not to the ownership of
the good.
\end{solution}

\begin{solution}{exo:b2:behavioural-evidence:5}
% code: exercises/solutions.py ex_prospect
$1000^{0.9} = 501.2$ and $400^{0.9} = 219.7$, so $V = 0.5 \times 501.2 - 0.5 \times 2 \times 219.7 = 250.6
- 219.7 = 30.9 > 0$: the gamble is accepted, narrowly. Its expected value is $0.5 \times 1000 - 0.5
\times 400 = \dal{300}$; expected utility with a nearly linear utility accepts it by a wide margin. A
gain two and a half times the loss is barely enough under \omterm{def:b2:behavioural-evidence:loss-aversion}{loss aversion} of two with diminishing
sensitivity.
\end{solution}

\begin{solution}{exo:b2:behavioural-evidence:6}
% code: exercises/solutions.py ex_default
% ledger: b2:ms2000-participation, b2:ms2000-default
Before: $0.37 \times 2000 = 740$ participants; after: $0.86 \times 2000 = 1720$, that is 980 more. At the
default: $0.61 \times 2000 = 1220$ employees, contributing $1220 \times 0.03 \times 30000 = \dal{1098000}$
a year. In the record, three quarters of the participants who chose for themselves contributed 6 per
cent or more; some of those who would have joined anyway and stay at the default contribute less than
they would have chosen, so the rise in participation overstates the rise in saving.
\end{solution}

\begin{solution}{exo:b2:behavioural-evidence:7}
% code: exercises/solutions.py ex_salience
Posted: $q = 1600 - 40 \times 22 = 720$. Since posting reduces demand by 8 per cent, unposted demand is
$720/0.92 = 782.61$. Then $1600 - 40(20 + 2\theta) = 782.61$, so $800 - 782.61 = 80\theta$ and $\theta =
17.39/80 = 0.217$: buyers act as if they saw about a fifth of the tax.
\end{solution}

\begin{solution}{exo:b2:behavioural-evidence:8}
% ledger: b2:camerer2016-abeler, b2:camerer2016-ericson
$0.08/0.18 = 0.44$ and $0.12/0.21 = 0.57$; the reported $0.43$ and $0.58$ use the unrounded effect
sizes. Accurate: ``Both replications found an effect in the original direction, about half as large as
the original, and not significant at the 5 per cent level.'' ``The evidence that expected earnings act
as a \omterm{def:b2:behavioural-evidence:reference-point}{reference point} is weaker than the original studies suggested.'' Overstated: ``The replications
showed that expectations do not affect \omterm{def:b2:behavioural-evidence:reference-point}{reference points}'' --- a non-significant estimate of half the
original size does not show an effect of zero.
\end{solution}

\begin{solution}{exo:b2:behavioural-evidence:9}
Give half the participants a good and run a market, as with the mugs; in a second treatment, with the
same institution, the same participants and the same procedures, trade tokens redeemable for a known
sum, with a different redemption value for each participant. Transaction costs (the effort of trading,
misunderstanding the rules) predict low volume in both treatments. \omterm{def:b2:behavioural-evidence:loss-aversion}{Loss aversion} over goods predicts low
volume for the good and the predicted volume for tokens, which are held only for exchange. Kahneman,
Knetsch and Thaler found the second pattern.
\end{solution}

\begin{solution}{exo:b2:behavioural-evidence:10}
% code: exercises/solutions.py ex_fourfold
(a) $\omega(0.01) = 0.047$, $\omega(0.05) = 0.122$, $\omega(0.5) = 0.439$, $\omega(0.9) = 0.745$,
$\omega(0.99) = 0.930$.
(b) The 1 per cent chance of \dal{1000} is worth $0.047 \times 1000 = 46.9$, more than the sure 10: risk
seeking for an unlikely gain. The 99 per cent chance is worth $930.3$, less than the sure 990: \omterm{def:b2:expected-utility:risk-aversion}{risk
aversion} for a likely gain.
(c) With losses, $\kappa$ multiplies both options and cancels. A 1 per cent chance of losing 1000 is
valued $-46.9\kappa$, worse than a sure loss of $10\kappa$: \omterm{def:b2:expected-utility:risk-aversion}{risk aversion} for an unlikely loss. A 99 per
cent chance of losing 1000 is valued $-930.3\kappa$, better than a sure loss of $990\kappa$: risk seeking
for a likely loss.
(d) \omterm{def:b2:expected-utility:risk-aversion}{Risk aversion} for likely gains and unlikely losses, risk seeking for unlikely gains and likely
losses. It is used to explain buying insurance against rare losses and \omterm{def:b2:expected-utility:lottery}{lottery} tickets at the same
time, and gambling to escape a loss that is almost certain.
\end{solution}

\begin{solution}{exo:b2:behavioural-evidence:11}
% code: exercises/solutions.py ex_eu_endowment
(a) WTP solves $\ln(m - P) + \gamma_m = \ln m$, so $m - P = me^{-\gamma_m}$ and $P = m(1 - e^{-\gamma_m})$.
WTA solves $\ln(m + A) = \ln m + \gamma_m$, so $A = m(e^{\gamma_m} - 1)$. The ratio is $(e^{\gamma_m} -
1)/(1 - e^{-\gamma_m}) = e^{\gamma_m}$.
(b) $e^{-\gamma_m} = 1 - 5/30000$, so $e^{\gamma_m} = 1/(1 - 1/6000) = 1.00017$: the owner asks
$\dal{5.0008}$, a gap of less than a tenth of a cent.
(c) Let $U(y, 1) = U(y, 0) + \Delta$, where $\Delta > 0$ is the value of the good, which by assumption
does not depend on $y$, and let $U_y$ be continuous and positive. WTP solves $U(m - P, 1) = U(m, 0)$, that
is $U(m, 0) - U(m - P, 0) = \Delta$; WTA solves $U(m + A, 0) - U(m, 0) = \Delta$. By the mean-value
theorem, $P\,U_y(m - s) = \Delta = A\,U_y(m + t)$ for some $s \in (0, P)$ and $t \in (0, A)$, so
\[
\frac{A}{P} = \frac{U_y(m - s)}{U_y(m + t)}.
\]
As $\Delta \to 0$, $P$ and $A$ tend to zero, hence $s$ and $t$ too, and by continuity the ratio tends to
$U_y(m)/U_y(m) = 1$. A ratio near two for a mug would need the marginal utility of money to halve over a
few dalens.
\end{solution}

\begin{solution}{exo:b2:behavioural-evidence:12}
% code: exercises/solutions.py ex_labour
(a) Below the target, an extra hour adds $\kappa w$ minus $\nu h$, so the interior choice is $h = \kappa
w/\nu = w/2$; above it, $h = w/\nu = w/4$. She stops at the target, $h = 200/w$, when the below-target
choice would overshoot it, $w \cdot w/2 \ge 200$, and the above-target choice would not reach it, $w \cdot
w/4 \le 200$: $20 \le w \le \sqrt{800} = 28.28$. Hours: at $w = 16$, $h = 8$ (income 128, below the
target); at $20$, $h = 10$; at $25$, $h = 8$; at $28$, $h = 200/28 = 7.14$; at $32$, $h = 8$ (income 256).
(b) Inside the band $h = 200/w$: elasticity $-1$. Outside, $h$ is proportional to $w$: elasticity $+1$.
(c) With $\kappa = 1$ there is no kink: $h = w/4$ always (4, 5, 6.25, 7, 8 hours), elasticity $+1$.
(d) Hours that rise with the wage on most driver-days, including days whose high earnings came as a
surprise, and quitting decisions that depend on the hours already worked rather than on the income
already earned.
\end{solution}

\begin{solution}{pb:b2:behavioural-evidence:1}
% code: exercises/solutions.py problem
% figdata: gamble_survey, gamble_summary (median_switch 200, share_reject_125 91.3)
\textbf{Part I.}
1. Essentially everyone: the gamble has expected value $\dal{12.5}$, and a smooth utility of wealth is
nearly linear over \dal{100}, so the \omterm{def:b2:expected-utility:risk-premium}{risk premium} is far smaller than 12.5.
2. $100 - 8.7 = 91.3$ per cent reject it.
3. Half accept from $z_+ = 200$. With linear value, $\kappa = 200/100 = 2.000$.
4. $200^{0.9} = \kappa \times 100^{0.9}$ gives $\kappa = 2^{0.9} = 1.866$.
5. Both lie in the interval from $1.820$ to $2.102$, on either side of the mean $1.955$
% ledger: b2:bivc2024
(Brown et al.\ 2024).
6. $(100 + 125) \times 125/(125 - 100) = 1125$: she must turn down every 50--50 gamble that loses
\dal{1125} or more, whatever it gains --- including a gamble that loses 1125 or wins a million.
\textbf{Part II.}
7. $8/3.5 = 2.286$.
8. Goods only: $\kappa = 2.286$. Money too: $\kappa = \sqrt{2.286} = 1.512$.
9. The goods-only version, $2.286$, is close to the survey's $2.000$; the other, $1.512$, is below the
survey and below the meta-analytic interval.
10. If tokens traded at about the predicted volume, the low trade in mugs is not due to the market's
costs or rules; if tokens also traded little, the gap may not be an \omterm{def:b2:behavioural-evidence:endowment-effect}{endowment effect} at all.
\textbf{Part III.}
11. $P_{\text{in}} = \Pr(b_i > \chi) = \Phi(\mu - \chi)$ and $P_{\text{out}} = \Pr(b_i > -\chi) =
\Phi(\mu + \chi)$.
12. $\mu - \chi = \Phi^{-1}(0.37) = -0.3319$ and $\mu + \chi = \Phi^{-1}(0.86) = 1.0803$. Then $\chi =
(1.0803 + 0.3319)/2 = 0.7061$ and $\mu = (1.0803 - 0.3319)/2 = 0.3742$.
13. $\mu_B - 0.7061 = \Phi^{-1}(0.45) = -0.1257$, so $\mu_B = 0.5804$.
14. $\Phi(0.5804 + 0.7061) = \Phi(1.2865) = 0.9009$: 90.1 per cent, a gap of $90.1 - 45.0 = 45.1$ points.
15. Opt-in: $0.45 \times 6 = 2.70$ per cent of pay per employee. Opt-out: $0.61 \times 3 + (0.901 - 0.61)
\times 6 = 1.83 + 1.75 = 3.58$ per cent. Opt-out raises saving here, but only because the extra
participants outweigh those held at 3 per cent.
16. $\Pr(-\chi < b_i < 0) = \Phi(1.2865) - \Phi(0.5804) = 0.9009 - 0.7192 = 0.1817$: 18.2 per cent.
17. $\Pr(0 < b_i < \chi) = \Phi(0.5804) - \Phi(-0.1257) = 0.7192 - 0.45 = 0.2692$: 26.9 per cent.
18. No. It counts people on the wrong side of zero, but to rank the defaults it would need the size of
their losses in dalens, whether $b_i$ is the employee's own welfare or her choice under \omterm{def:b2:intertemporal-choice:present-bias}{present bias},
whether $\chi$ is a real cost or an error, and the cost of the employer's matching contribution.
\textbf{Part IV.}
19. $0.05 \times 500 = \dal{25}$ a year, against \dal{100} of extra premium. The extra risk of the high
\omterm{def:b2:insurance-and-risk-sharing:deductible}{deductible} has a variance of $0.05 \times 0.95 \times 500^2 = 11875$; with a relative \omterm{def:b2:expected-utility:risk-aversion}{risk aversion} of
10 at a wealth of \dal{50000}, its \omterm{def:b2:expected-utility:risk-premium}{risk premium} is about $\tfrac12 \times (10/50000) \times 11875 = 1.19$
dalens. Expected utility predicts the high \omterm{def:b2:insurance-and-risk-sharing:deductible}{deductible}.
20. $\omega(0.05) = 0.1215$. The low \omterm{def:b2:insurance-and-risk-sharing:deductible}{deductible} is bought when $\kappa \times 0.1215 \times 500 > 100$,
that is when $\kappa > 100/60.75 = 1.646$.
21. $\Pr(\kappa > 1.646) = 1 - \Phi\bigl((\ln 1.646 - \ln 2)/0.35\bigr) = 1 - \Phi(-0.5566) = 0.711$:
71.1 per cent.
22. The insurer's records of the residents' \omterm{def:b2:insurance-and-risk-sharing:deductible}{deductible} choices, linked to their survey answers. The model
is rejected if far fewer than about 70 per cent choose the low \omterm{def:b2:insurance-and-risk-sharing:deductible}{deductible}, or if the choice is unrelated
to each resident's measured $\kappa$.
23. A misperceived (overestimated) chance of a claim; difficulty in paying \dal{1000} at once; the
insurer's choice of a default \omterm{def:b2:insurance-and-risk-sharing:deductible}{deductible} or its sales advice; and the premium being coded as a loss too,
which would work against the low \omterm{def:b2:insurance-and-risk-sharing:deductible}{deductible}.
24. The survey gives $\kappa = 2.000$ with linear value ($1.866$ with curvature $0.9$), inside the
meta-analytic interval; the default model predicts 90.1 per cent participation under opt-out against
45 per cent under opt-in, a gap of 45.1 points; and \omterm{def:b2:behavioural-evidence:prospect-theory}{prospect theory} predicts that 71.1 per cent of
residents buy the low \omterm{def:b2:insurance-and-risk-sharing:deductible}{deductible}, which expected utility says almost no one should.
\end{solution}
